\abstract{Disease early warning and surveillance (DEWS) systems in low- and middle-income countries (LMICs) face persistent architectural failures: data centralization, metadata-level privacy exposure, unscalable ledger replication, absence of selective disclosure, and exclusion of resource-constrained peripheral sites from meaningful system participation. Existing blockchain-health architectures exacerbate these weaknesses by treating the ledger as a bulk data repository, assuming uniform node capacity, and conflating immutability with privacy. This paper presents a hybrid, privacy-layered surveillance architecture that replaces full on-chain storage with selective anchoring: encrypted surveillance payloads reside off-chain via IPFS, while Ethereum smart contracts retain only cryptographic hashes, access-control state, provenance records, alerts, and audit events. The architecture introduces five integrated novelty contributions: (1)~hierarchical event processing with district-level batching aligned to public-health administrative tiers; (2)~a layered privacy model combining zero-knowledge verification, dynamic grant/revoke/expire permissions, and encrypted off-chain records; (3)~edge and fog preprocessing to reduce bandwidth, latency, and compute burden at peripheral facilities; (4)~a comparative consensus framework evaluating Proof of Authority, Delegated Proof of Stake, PBFT variants, and Raft-based alliance models against LMIC workload, trust, and infrastructure constraints; and (5)~FHIR-aware interoperability workflows supporting standards-compliant selective disclosure across organizational boundaries. A functional prototype validates the core smart-contract coordination layer. Performance evaluation demonstrates stable block generation (12--14\,s), usability assessment with 26 domain professionals yields Software Usability Measurement Inventory scores significantly exceeding standardized benchmarks, and an expanded threat model addresses validator collusion, metadata leakage, overbroad tier exposure, and equity-aware governance. The architecture advances the state of the art by treating privacy, consensus, and resource asymmetry as linked governance concerns rather than isolated engineering parameters, yielding a scalable, privacy-preserving, and LMIC-optimized surveillance infrastructure.}

\keywords{blockchain, disease early warning, surveillance, Ethereum, smart contracts, health informatics, distributed systems}

\jnlcitation{\cmark[Authors]}
